\chapter[Orthogonality Is Not Enough]{Orthogonality Is Not Enough: A Critique of Double Machine Learning}\label{ch:critique}

\begin{tcolorbox}[colback=gray!6,colframe=gray!50,title=In brief,fonttitle=\normalfont\itshape,coltitle=black,colbacktitle=gray!15,boxrule=0.4pt,arc=1pt]
Orthogonality removes the first order effect of a nuisance error. What it leaves is a signed difference of two channels, and the standard rate condition bounds each in absolute value, so it cannot see when they cancel. The reported interval comes from one nuisance fit and omits how the bias varies across fits: cutting the folds from five to two raises the spread of the estimator by 12 to 72 percent in the sparse designs while the reported standard error moves by about five percent at most. The partially linear score targets a variance weighted average, and no learner gives an interval of finite expected length that is valid uniformly in overlap. A posterior over the nuisance pair integrates over the fit and reports its own spread, but a Bayesian analysis with the prior on the structural coefficients does far worse than DML. The evidence is mostly simulated, in linear Gaussian designs with $n = 300$ and $p = 400$.
\end{tcolorbox}

\section{What the critique is and is not}\label{sec:crit:what}

Take one data set from the weak overlap design\index{simulation study!weak overlap, critique|textbf} of Polson and Sokolov (2026b): $n = 300$ observations, $p = 400$ controls of which fifteen confound, a treatment residual with standard deviation 0.3, and a true effect $\tau = 1$. Fit double machine learning\index{double machine learning} with a cross validated lasso at five folds, then fit it 199 more times, changing only the random assignment of observations to folds. The 200 estimates average 1.427 with a standard deviation of 0.075, and the reported standard error is 0.196. On 141 of the 200 assignments the 95 percent interval excludes the truth.

The assignments are not what puts the truth outside the interval: the spread of 0.075 sits well inside the reported error of 0.196. The estimates average 2.2 reported standard errors above $\tau$, so the interval has about the right width around the wrong point, as the paper reports. A posterior from the reduced form Bayesian model of Section~\ref{sec:crit:bayes} on the same data has mean 1.208, standard deviation 0.249 and 95 percent interval 0.722 to 1.692. It covers the truth and is not centred on it either.

DML replaces a distribution over nuisance functions by a point estimate, then certifies the consequences using only that estimate\index{double machine learning!critique of|textbf}. The certificate needs functionals of the law of the nuisance error, which a single fit cannot supply. Chapter~\ref{ch:dml} states the method; this chapter follows Polson and Sokolov (2026b).

DML is not wrong. It is correct under its assumptions, and its uniformity statement over a class defined by explicit constants is a theorem (Chernozhukov and coauthors 2018). The claim is that membership in the class cannot be certified from the data, that the class excludes regimes reached at ordinary sample sizes, and that the reported interval omits a variance component it cannot report. Nor does Bayes win by default: Section~\ref{sec:crit:bayes} exhibits a Bayesian analysis that does far worse than DML, and the comparison that remains is mixed.

Unless stated otherwise, simulations use the linear Gaussian baseline of Section~\ref{sec:dml:practice}, with $n = 300$, $p = 400$ and outcome noise variance 0.9. Tables~\ref{tab:crit:folds}, \ref{tab:crit:head} and~\ref{tab:crit:cmp} and their arm comparisons use our rerun of the Chapter~\ref{ch:dml} data sets: 5,000 replications per design at $s = 5$, 15 and 30 and under weak overlap, every arm on the same data (Appendix~\ref{app:code}). The dense rows of Table~\ref{tab:crit:head}, the channel measurements (5,000 replications per design), the cross fitted posterior and the timings are also ours, with seeds shared across arms. Apart from these and the other runs of our own that the text names, simulated DML numbers and the nine procedure comparisons on the three real data sets of Section~\ref{sec:crit:real} are quoted from Polson and Sokolov (2026b). The paper uses the cross fitted partially linear score with a cross validated lasso or, where stated, a horseshoe posterior mean, another nuisance learner or the augmented score, with 500 replications per cell and shared seeds. Its lasso and horseshoe baselines differ from our rerun's by up to three points of coverage. The dense design gives every coordinate of both coefficient vectors a nonzero value, spreading the total signal of the $s = 15$ design over all $p$ coordinates so that no sparse learner can find it. Two further runs of our own use 1,000 data sets per design: the refitting bootstrap (Section~\ref{sec:crit:se}) and the random forest, gradient boosting and stack comparison (Section~\ref{sec:crit:learner}); Appendix~\ref{app:code} names their scripts. Neural networks are never run. Table~\ref{tab:dml:compare} contrasts the approaches by construction.

\begin{table}[ht]
\centering
\small
\caption{What double machine learning and a Bayesian causal model deliver by construction. The entries describe what each is built to provide, not which does better in a given data set.}
\label{tab:dml:compare}
\begin{tabular}{@{}L{1.15in}L{1.5in}L{1.5in}@{}}
\toprule
 & DML & Bayesian model \\
\midrule
Average effect & efficient under rate conditions with the augmented score & efficient when a semiparametric Bernstein von Mises theorem holds, shown for Gaussian process priors with a propensity component and not for the horseshoe \\
Uncertainty & interval from one nuisance fit & posterior over the nuisance pair, calibrated only if the model is right \\
Conditional effects & R learner or causal forest & posterior for $\tau(x)$ \\
Other functionals & a new score for each & draws of the same posterior \\
Misspecification & robust to one nuisance error at a time with the augmented score & depends on the flexibility of the prior \\
Decisions & not directly & posterior expected utility \\
Partial identification & not addressed & prior over the unidentified part \\
Computation & a few regressions & sampling or amortized training \\
\bottomrule
\end{tabular}
\end{table}

\section{The error in two channels}\label{sec:crit:channels}

 By orthogonality\index{Neyman orthogonality!is not enough|textbf} the first order effect of the two fitting errors is gone. What remains can be written down exactly.

Write $\delta_m = \hat m - m$ and $\delta_\ell = \hat\ell - \ell$ for the treatment and outcome fitting errors, and let $\|\cdot\|$ and $\langle\cdot,\cdot\rangle$ be the norm and inner product in $L_2$ of the covariate law. Lemma~\ref{dml:lem:exact} writes the error of the cross fitted estimator as $(S_n + C_n + B_n)/\hat J_n$, with $\hat J_n$ the mean squared treatment residual, and Corollary~\ref{dml:cor:channels} gives the mean of $B_n$ as the cross channel $\langle\delta_m,\delta_\ell\rangle$ minus the own square channel\index{error decomposition!two channels} $\tau\|\delta_m\|^2$, so that at the realized fits the leading bias is
\begin{equation}
\Delta = \frac{b}{J + \|\delta_m\|^2}, \qquad b = \langle\delta_m,\delta_\ell\rangle - \tau\|\delta_m\|^2,
\label{eq:crit:bias}
\end{equation}
with $J = \E\,\mathrm{Var}(T \mid X)$ measuring overlap\index{overlap}.

The usual argument bounds the cross channel by Cauchy and Schwarz, $|\langle\delta_m,\delta_\ell\rangle| \le \|\delta_m\|\,\|\delta_\ell\|$, adds the own square channel, and asks that $\|\delta_m\|(\|\delta_m\| + \|\delta_\ell\|) = o_p(n^{-1/2})$, the rate condition of Theorem~\ref{dml:thm:main}. The bound on the cross channel is tight only when the two error functions point the same way. Write $\rho_\delta = \langle\delta_m,\delta_\ell\rangle/(\|\delta_m\|\,\|\delta_\ell\|)$, when both norms are positive, for the cosine between them, the alignment coefficient\index{alignment coefficient|textbf}, so that the cross channel is exactly $\rho_\delta\|\delta_m\|\,\|\delta_\ell\|$ and the rate condition\index{rate condition} is a worst case over $\rho_\delta$. Two features of \eqref{eq:crit:bias} matter. The channels enter with opposite signs, so a learner that shrinks one can enlarge the total. And because the rate condition bounds each channel in absolute value, it cannot credit their cancellation: it asks for $\|\delta_m\| = o_p(n^{-1/4})$ even when the two channels nearly cancel. The condition that governs the estimator is the one in the next theorem\index{error decomposition!sharp bound on the bias|textbf}.

\begin{theorem}[Sharp rate condition]\label{crit:thm:sharp}
Let $\E(\nu \mid X) = 0$, $\E(\varepsilon \mid T, X) = 0$, $\mathrm{Var}(\nu \mid X) = \sigma_\nu^2 > 0$, $\E(\varepsilon^2 \mid \nu, X) = \sigma_\varepsilon^2$ and $\E\varepsilon^4 + \E\nu^4 < \infty$. Let the fits be cross fitted with $K$ fixed, $\|\delta_m\| + \|\delta_\ell\| = o_p(1)$ for the fits of every fold, and $\E\delta_m^4 \le C\|\delta_m\|^4$, $\E\delta_\ell^4 \le C\|\delta_\ell\|^4$ for a constant $C$. Here $b$ is the fold weighted average $\sum_k (n_k/n)\,b^{(k)}$, where $b^{(k)}$ is the value of \eqref{eq:crit:bias} at the fits trained without fold $k$. If $\sqrt n\,b = o_p(1)$\index{error decomposition!sharp bound on the bias}, then $\sqrt n(\hat\tau - \tau)$ converges to a normal law with mean zero and variance $\sigma_\varepsilon^2/\sigma_\nu^2$. Proof in Appendix~\ref{app:dmlproofs}.
\end{theorem}

The condition $\sqrt n\,b = o_p(1)$ is implied by the rate condition\index{rate condition}, since $|b| \le \max(1, |\tau|)\,\|\delta_m\|(\|\delta_m\| + \|\delta_\ell\|)$, and is strictly weaker: it also holds when the cross channel is small because $\rho_\delta$ is small and not because $\|\delta_m\|\,\|\delta_\ell\|$ is, and when $\rho_\delta \to 1$ but the two channels cancel. The converse of the theorem holds under two further hypotheses, asymptotic independence of $\sqrt n(b + C_n)$ and $\sqrt n S_n$ and uniform integrability, which the paper assumes and does not derive. One consequence concerns diagnostics. Since $\hat J_n$ converges to $J + \|\delta_m\|^2$, the denominator DML reports overstates overlap by the treatment fitting error, and it should not be read as an overlap diagnostic.

Cancellation deserves a look of its own, because it makes the bias small without making either channel small. Suppose both regressions are sparse on a common support of size $s$ with agreeing signs and are fitted by the lasso on an orthogonal design with one common penalty $\lambda_{\mathrm L}$. Both channels then equal $\lambda_{\mathrm L}^2 s$ to leading order, and $b = (1 - \tau)\lambda_{\mathrm L}^2 s$ vanishes at $\tau = 1$ (Exercise~\ref{ex:chcrit:6}). The cancellation rests on the common penalty. Cross validation picks a penalty for each equation, and in our designs the outcome equation has the larger noise (variance 1.9 against 1), so the penalties differ and $b = \lambda_{\mathrm L}^m s(\lambda_{\mathrm L}^\ell - \tau\lambda_{\mathrm L}^m)$ is positive at $\tau = 1$ and grows with $s$. At $\tau = 1$ and two folds the two channels stay close and their difference is ten to thirteen percent of the cross channel: 0.30 against 0.27 at $s = 5$ and 3.72 against 3.24 at $s = 30$, with biases of 0.023 and 0.106. Moving one thing at a time shows the cancellation is near and not exact. With 15 true confounders and five folds, $\tau = 0.5$ gives channels 0.302 and 0.199 and coverage 0.685, $\tau = 1$ gives 0.451 and 0.401 and coverage 0.826, and $\tau = 2$ gives 0.814 and 0.801 and coverage 0.864. A smaller effect makes DML worse and a larger one better, so $\tau = 1$ is not the favourable end of the range.

\section{When the bound is attained}\label{sec:crit:align}

Suppose both regressions are sparse and share their active covariates, which is what confounding usually is. The next proposition says what a lasso does to the two fitting errors then. Its hypotheses are ones that cross validation does not meet.

\begin{proposition}[Alignment under retained shrinkage]\label{crit:prop:alignment}
Suppose the treatment and outcome regressions are linear and sparse on a common support $S$ of size $s$ with agreeing signs, the covariates are sub-Gaussian with identity covariance, both regressions are fitted by the lasso with one common penalty $\lambda_{\mathrm L} \asymp \sigma\sqrt{\log p/n}$ and the shrunken fits are kept, and the lasso recovers $S$ exactly (which needs a beta min and an irrepresentability condition), with $p \to \infty$ and $s\log^2 p = o(n)$. Then on the event of exact recovery $\langle\delta_m,\delta_\ell\rangle = \lambda_{\mathrm L}^2 s\,\{1 + o_p(1)\}$ and $\|\delta_m\|^2 = \|\delta_\ell\|^2 = \lambda_{\mathrm L}^2 s\,\{1 + o_p(1)\}$. Hence $\rho_\delta \to 1$ in probability\index{alignment coefficient!when the bound is attained}, and $\sqrt n\,\langle\delta_m,\delta_\ell\rangle \to 0$ if and only if $s = o(\sqrt n/\log p)$. Proof in Appendix~\ref{app:dmlproofs}.
\end{proposition}

The lasso shrinks every retained coefficient toward zero by the same amount. When the same covariates matter in both equations with the same signs, the two fits are pulled the same way by the same distance, and each error is mostly a copy of the other. The Cauchy and Schwarz bound inside the rate condition is then attained\index{alignment coefficient!when the bound is attained|textbf}, and it forces $s = o(\sqrt n/\log p)$. A least squares refit on the recovered support removes the shrinkage and does not pay this price. Every learner that shrinks strong coefficients by a fixed amount and reports the shrunken fit does.

The hypotheses fail in practice, and the data say so. Cross validation chooses a penalty for each equation and essentially never recovers the support: the exact recovery frequency is at most 0.001 at $s = 5$, 15 and 30 over 5,000 replications, and at five folds the selected sets hold 26.1, 64.3 and 112.9 variables against true supports of 5, 15 and 30, all of the error being false positives. The one run in which exact recovery holds uses DML with a lasso at the theory penalty of Belloni, Chernozhukov and Hansen (each equation's penalty scaled to its own noise level, shrunken fits kept) at $s = 5$ and five folds. In the treatment equation it selects the true support and nothing else in 99.9 percent of 25,000 fits (five folds in each of 5,000 replications). There the alignment coefficient\index{alignment coefficient} rises from 0.83 to 0.99, the bias from 0.017 to 0.177, and coverage falls from 0.917 to 0.034. The two penalties differ, so the cancellation of the common penalty case does not apply and the bias is the positive one of the separate penalty formula above. The estimator that recovers the support is the worse one by a wide margin.

What the cross validated lasso does show is a measured alignment that is high and flat. At two folds $\rho_\delta$ is 0.815, 0.812 and 0.815 at $s = 5$, 15 and 30, with a Monte Carlo standard error of 0.002, and the horseshoe nuisance of Section~\ref{sec:crit:learner} sits between 0.48 and 0.66. The measurement has limits. It is a two fold number with separate cross validated penalties, in designs that share support and signs by construction and have $\tau = 1$ and identity covariance. The sign mechanism of the proposition explains little of it: making the two coefficient vectors far less collinear, with the cosine falling from 0.999 to 0.53 and the signs still agreeing, moves $\rho_\delta$ by only 0.03 to 0.04. We attribute the rest to the two fits selecting largely the same false positives, so that shared selection noise is a common part of both errors. That reading has no theory behind it. And $\rho_\delta$ needs the truth, so the analyst cannot compute it.

\section{What a better learner buys}\label{sec:crit:learner}

The natural response to a shrinkage bias\index{shrinkage bias|textbf} that does not decay with the size of the signal is a learner\index{nuisance function!learners for|textbf} whose bias does\index{error decomposition!what a better learner buys|textbf}. The horseshoe posterior mean (Chapter~\ref{ch:horseshoe}) shrinks small coefficients hard and leaves strong ones with a bias of order $1/|y|$. In the orthogonal Gaussian sequence model it does what is hoped for the cross channel\index{horseshoe prior!oracle property|textbf}.

\begin{remark}[The horseshoe cross channel, in the sequence model]\label{crit:thm:hsoracle}
Polson and Sokolov (2026b) prove the following in the orthogonal Gaussian sequence model, on the signal set $S$ only. Let $s = |S| \to \infty$, let the horseshoe global scale $\gscale$ be estimated by marginal likelihood truncated to $[1/p, 1]$, so that $\gscale \lesssim (s/p)\sqrt{\log(p/s)}$ with probability tending to one, and let $\tau \ne 0$. Write $\sigma_\ell^2 = \tau^2\sigma_\nu^2 + \sigma_\varepsilon^2$ and $\sigma = \max(\sigma_\nu, \sigma_\ell)$, and assume the beta min condition $\min_{j \in S}(|\gamma_j| \wedge |\varsigma_j|) \gg \sigma\sqrt{\log(p/s)/n}$ and the noise condition $\sigma_\ell = o\{\sigma_\nu\sqrt{s\log(p/s)}\}$, which holds automatically when the noise scales are fixed. Then the horseshoe cross channel on the signal set is $\tau\sigma_\nu^2 s/n\,\{1 + o_p(1)\}$, the value for an oracle that knows the support and runs least squares on it\index{horseshoe prior!oracle property}. For the lasso\index{lasso} the same channel is $\lambda_{\mathrm L}^2 s \asymp 2\sigma^2 s\log p/n$, larger by a factor of order $\log p$, so the cross channel alone would allow $s = o(\sqrt n)$ for the horseshoe against $s = o(\sqrt n/\log p)$ for the lasso. The proof is in the paper and is not reproduced here. The extension to a general design is open.
\end{remark}

Two cautions come with it. The sequence model guides the cross channel and does not model a regression with $p > n$. And the size of the effect is off. The leading constant of the ratio is $2\sigma_\ell^2\log p/(|\tau|\sigma_\nu^2)$, with $\sigma_\ell^2 = \tau^2\sigma_\nu^2 + \sigma_\varepsilon^2 = 1.9$ in our design, which is 22.8 at $p = 400$. The measured ratios of the lasso's cross channel to the horseshoe's at two folds are 1.2, 1.9 and 4.0 at $s = 5$, 15 and 30. The direction and the growth with $s$ are supported, and the constant is overpredicted by factors of 19, 12 and 6.

The own square channel is the other half, and the horseshoe improves it by less.

\begin{remark}[The own square channel keeps its logarithm in the worst case]\label{crit:prop:own}
Over treatment coefficient vectors with $s$ nonzero entries, the worst case squared error of the lasso is of order $\sigma^2 s\log p/n$, by the prediction bound of Bickel, Ritov and Tsybakov (2009) and a matching lower bound from Proposition~\ref{crit:prop:alignment}, while in the sequence model the horseshoe attains the minimax rate $\sigma_\nu^2 s\log(p/s)/n$ (van der Pas, Kleijn and van der Vaart 2014). When $s \le p^a$ for a fixed $a < 1$, so that $\log(p/s) \asymp \log p$, the two worst cases agree up to constants, and no factor $\log p$ separates them (Appendix~\ref{app:guarantees} states the same qualification for adaptive contraction).
\end{remark}

This is a worst case over the sparsity class, attained near the detection boundary, and it is not the pointwise comparison. For signals well above the boundary the horseshoe's risk is of order $\sigma_\nu^2 s/n$ and it improves the own square channel too, so the logarithm survives only in the worst case. The measurements say by how much. Going from the lasso to the horseshoe at two folds divides the own square channel by 0.7, 1.4 and 3.1 at $s = 5$, 15 and 30, and the cross channel by 1.2, 1.9 and 4.0. The relief on the cross channel is larger at every $s$, and at $s = 5$ the own square channel gets worse, by about 40 percent, from 0.27 to 0.39.

Because the bias is a difference, that ordering decides its sign. At $s = 5$ the horseshoe has the smaller cross channel, 0.25 against 0.30, and the larger bias with the opposite sign. The numerators of the leading bias are $0.247 - 0.385 = -0.138$ for the horseshoe and $0.301 - 0.273 = +0.028$ for the lasso, and the measured biases are $-0.084$ and $+0.024$. The fold count then moves both channels. Run at five folds, so that each fit sees 240 observations and not 150, the same horseshoe nuisance on the same data sets has channels 0.045 and 0.055 in place of 0.247 and 0.385, and a bias of $-0.009$ in place of $-0.084$, while its alignment coefficient moves the other way, from 0.48 to 0.61. A worse cosine came with a bias an order of magnitude smaller. The mixed sign design says the same from the other side. In our rerun at five folds (5,000 replications), alternating the signs of the direct outcome coefficients along the support, with the same magnitudes, lowers the cosine between the two coefficient vectors from 0.999 to 0.734 and the lasso's $\rho_\delta$ at $s = 15$ from 0.815 to 0.571, and the bias does not fall with it. At $s = 15$ it changes sign and grows from $+0.037$ to $-0.100$, and at $s = 30$ it reaches $-0.194$ with coverage 0.106 against 0.734. At $s = 15$ the cross channel is now 0.259 against an own square channel of 0.401, and $\rho_\delta$ does not enter the second. The angle between the errors does not decide the bias. The difference of the channels does.

When the cross channel is small beside the own square channel, DML turns into a familiar problem.

\begin{theorem}[Attenuation representation]\label{crit:thm:atten}
Assume the conditions of Theorem~\ref{crit:thm:sharp} other than $\sqrt n\,b = o_p(1)$, and write $Q_n = \sum_k (n_k/n)\|\delta_m^{(k)}\|^2$ and $H_n = \sum_k (n_k/n)\langle\delta_m^{(k)},\delta_\ell^{(k)}\rangle$ for the fold weighted channels, with $\delta_m^{(k)}$ and $\delta_\ell^{(k)}$ the errors of the fits trained without fold $k$. Suppose that $H_n = o_p(Q_n)$ and that $\sqrt n\,Q_n \to \infty$ in probability. Then at the realized fits
\[
\hat\tau - \tau = -\tau(1 - \mathcal R) + o_p(1 - \mathcal R), \qquad \mathcal R = \frac{J}{J + Q_n},
\]
so that $\hat\tau \approx \tau\mathcal R$\index{shrinkage bias!attenuation theorem|textbf}. Proof in Appendix~\ref{app:dmlproofs}.
\end{theorem}

The fitted treatment residual is the true residual observed with error. DML is least squares with a mismeasured regressor, and $\mathcal R$ is the classical reliability ratio\index{reliability ratio|textbf}: the estimate is multiplied by the share of the variance of the fitted residual that is signal. Nothing in the paper delivers the hypothesis, and the cross channel is not negligible: its measured ratio to the own square channel at two folds is 0.64, 0.81 and 0.88 for the horseshoe and 1.10, 1.12 and 1.15 for the lasso. Since $\|\delta_m\| \to 0$ sends $\mathcal R$ to one, the attenuation is a finite sample effect, and the representation is an approximation for samples in which the treatment fitting error still dominates the root $n$ noise. We use it for the mechanism and for the sign of the horseshoe's bias, which is negative at every $s$, not for the size of the bias. Section~\ref{sec:crit:real} meets an instance.

Forests and boosting do not escape the mechanism, and in this law they are far worse than the lasso. In a run of our own we put five learners in both nuisance regressions of the cross fitted partially linear score, each tuned only inside its training fold: the cross validated lasso; a random forest of 100 trees whose share of covariates tried at a split, 0.1 or 0.33, is chosen by out of bag fit; gradient boosted trees of depth 3 with early stopping; the horseshoe posterior mean from 500 Gibbs sweeps; and a stack that weights the out of fold predictions of the lasso, the forest and the boosting by non negative least squares. The law is the sparse partially linear one of this chapter with $p = 400$ and $\tau = 1$, at $s = 5$, 15 and 30, treatment residual standard deviation $\sigma_\nu = 1$ and 0.3, and $n = 300$ and 1,200, with five folds and 1,000 data sets per cell, every learner on the same data sets (Figure~\ref{fig:crit:learners}). The forest and the boosting have biases from 0.13 to 0.98 against $\tau = 1$, alignment coefficients from 0.78 to 0.99, and intervals that cover in at most 2 percent of the data sets in any cell. Their $\sqrt n\,b$ grows from $n = 300$ to 1,200 in every cell but one (boosting at $s = 5$ and $\sigma_\nu = 1$), by a factor of 1.6 to 2.0 for the forest, so their bias does not fall at the root $n$ rate. The horseshoe nuisance and the stack match the lasso where the lasso does well and beat it where it does not. At $\sigma_\nu = 1$ and $n = 1{,}200$ all three cover between 0.919 and 0.945 for $s \le 15$, but at $s = 30$ the lasso covers 0.884 against 0.930 and 0.921, and under weak overlap at $n = 300$ and $s = 15$ it covers 0.561 against 0.905 and 0.915. Their alignment under weak overlap is 0.22 to 0.35, against 0.57 to 0.72 for the lasso, which is the mechanism of Section~\ref{sec:crit:align} seen from the other side. The law is linear and Gaussian, which favours the linear learners, so this shows that a nonlinear learner is no escape here, not that it fails everywhere.

\begin{figure}[ht]
\centering
\includegraphics[width=\linewidth]{critique_learners.pdf}
\caption{Coverage of the nominal 95\% interval (top) and signed bias (bottom) of the cross fitted partially linear estimate at $n = 300$ and 1,200, by nuisance learner, with 95\% Monte Carlo intervals. Columns are $s = 5$, 15 and 30 at $\sigma_\nu = 1$ and then at $\sigma_\nu = 0.3$; $p = 400$, $\tau = 1$, five folds, 1,000 data sets per point, the same data sets for every learner.}
\label{fig:crit:learners}
\end{figure}

\section{What the reported standard error leaves out}\label{sec:crit:se}

Table~\ref{tab:crit:folds} compares five and two folds\index{cross fitting!number of folds|textbf} for the cross validated lasso on the same 5,000 data sets, with shared seeds and nothing else changed.

\begin{table}[ht]
\centering
\small
\caption{The fold count alone. DML with a cross validated lasso, $n = 300$, $p = 400$, $\tau = 1$, the same 5,000 data sets at five and at two folds, $s$ true confounders. Coverage is for nominal 95 percent intervals. The last two columns are the mean reported standard error and the Monte Carlo standard deviation of $\hat\tau$.}
\label{tab:crit:folds}
\begin{tabular}{@{}llrrrr@{}}
\toprule
$s$ & folds & coverage & bias & reported s.e. & s.d. of $\hat\tau$ \\
\midrule
5 & 5 & 0.917 & 0.017 & 0.0533 & 0.0586 \\
 & 2 & 0.865 & 0.023 & 0.0529 & 0.0654 \\
15 & 5 & 0.826 & 0.037 & 0.0523 & 0.0644 \\
 & 2 & 0.694 & 0.053 & 0.0513 & 0.0786 \\
30 & 5 & 0.734 & 0.050 & 0.0513 & 0.0734 \\
 & 2 & 0.450 & 0.106 & 0.0487 & 0.1265 \\
\bottomrule
\end{tabular}
\end{table}

Coverage\index{coverage}\index{coverage!of DML intervals} falls, which is no surprise: with fewer folds each nuisance is fitted on less data. McNemar's exact test gives $p = 3\times 10^{-30}$, $5\times 10^{-78}$ and $10^{-213}$. The columns to read are the last two. The reported standard error\index{standard error|textbf} falls by 0.7, 1.8 and 5.2 percent, while the Monte Carlo standard deviation of the estimator rises by 12, 22 and 72 percent. The ratio of reported error to spread goes from 0.91 to 0.81 at $s = 5$ and from 0.70 to 0.38 at $s = 30$, where the estimator becomes nearly three quarters again as variable and its interval gets narrower. The horseshoe nuisance shows the same thing: coverage falls from 0.928 to 0.577, from 0.895 to 0.587 and from 0.827 to 0.671 at $s = 5$, 15 and 30, and from 0.908 to 0.841 under weak overlap, as its spread rises by 142, 88, 24 and 27 percent against rises of 15, 9, 3 and 7 percent in its reported error. These are findings about two learners at $n = 300$ and $p = 400$ in linear designs. The change also cuts the training share of each fit from 80 to 50 percent, so the rise in spread mixes fold randomness with worse fits (the horseshoe arm's channels fall by factors of 5.5 and 7 when its fits use 240 observations and not 150). The cleaner statement concerns the reported standard error: when $K$ changes it moves by at most 5.2 percent for the lasso and 15 percent for the horseshoe, far less than the spread.

The lottery isolates the randomization\index{cross fitting!variance from the random split|textbf}. Hold one data set fixed ($s = 15$) and redraw only the fold assignment, 200 times at each of $K = 2$, 5 and 10. At $K = 5$ the estimates range from 0.932 to 1.069. Their standard deviation across assignments is 0.057, 0.026 and 0.019, against a reported standard error of 0.053 at all three fold counts. At five folds the randomization alone is half of what is reported, and at two folds it exceeds all of it. The horseshoe nuisance is no better: 0.102 at two folds against a reported 0.060, and 0.025 at five against 0.054. This data set's estimate happens to sit near the truth, so no five fold assignment rejects $\tau = 1$, and the number to read is the spread.

Real data show the same lottery. In the 401(k) data of Chapter~\ref{ch:k401}, repeating the five fold split gives lasso estimates from 9.52 to 10.26 thousand dollars over 20 partitions and forest estimates from 8.59 to 9.05 over 5, and Figure~\ref{fig:dmlsplits} draws every partition with its interval. The intervals are wide relative to the spread, so with the full sample the point estimate depends on a random seed without the conclusion changing. In our own extract of those data (a different set of 99 controls) the split to split standard deviation is a sixth of the reported standard error at $n = 9{,}915$ and two fifths in a random subsample of $n = 300$, where 200 fold assignments of that one data set return estimates between 1.90 and 26.22 thousand dollars.

\begin{figure}[ht]
\centering
\includegraphics[width=0.9\linewidth]{dml_splits.pdf}
\caption{Cross fitted DML estimates of the effect of 401(k) eligibility on net financial assets, with 95\% intervals, for 20 random fold partitions with lasso nuisances and 5 with random forest nuisances.}
\label{fig:dmlsplits}
\end{figure}

Why the reported error barely responds to $K$ is a statement about what it conditions on. Fix fits $\eta = (\tilde m, \tilde\ell)$ and write $h_\eta = (T - \tilde m(X))(Y - \tilde\ell(X))$, $q_\eta = (T - \tilde m(X))^2$, $\tau_\eta = \E h_\eta/\E q_\eta$ and $\psi_\eta = (h_\eta - \tau_\eta q_\eta)/\E q_\eta$, the influence function, with expectations and variances below over one fresh observation. A direct calculation gives $\tau_\eta = \tau + \Delta(\eta)$, with $\Delta$ the bias of \eqref{eq:crit:bias} at the errors of $\eta$: $\Delta$ is the bias a fit carries and $\psi_\eta$ the noise around it. The cleanest statement is for a single split\index{cross fitting!single split|textbf}.

\begin{theorem}[What the sandwich sees, and what it cannot]\label{crit:thm:single}
Split the sample into a training part and an evaluation part of size $n_E \ge 17$\index{cross fitting!single split}, fit $\hat\eta = (\hat m, \hat\ell)$ on the training part alone, and compute the Robinson estimator $\hat\tau$ on the evaluation part. Assume the moment conditions and fourth moment bounds of Theorem~\ref{crit:thm:sharp} (and, for the bounds on $\tilde\delta_{n_E}$ and $R_{n_E}$ only, finite eighth moments of $\nu$ and $\varepsilon$ and $\E\delta_m^8 \le C\|\delta_m\|^8$, $\E\delta_\ell^8 \le C\|\delta_\ell\|^8$), that $\nu$ given $X$ has a bounded density, and that the fits satisfy $\|\delta_m\| + \|\delta_\ell\| \le 1$. Then $\E\hat\tau^2 < \infty$ and, conditionally on the training part,
\[
\E(\hat\tau \mid \hat\eta) = \tau + \Delta(\hat\eta) + \delta_{n_E}, \qquad \mathrm{Var}(\hat\tau \mid \hat\eta) = \frac{\mathrm{Var}(\psi_{\hat\eta})}{n_E} + \tilde\delta_{n_E},
\]
where $\|\delta_{n_E}\|_2 \le C_1/n_E$ and $\|\tilde\delta_{n_E}\|_1 \le C_2\,n_E^{-3/2}$, with constants that depend only on $J$, the moment bounds and the density bound. The sandwich variance estimator, which recentres at $\hat\tau$, estimates the leading term of the second display. Unconditionally, with $\sigma_\Delta^2 = \mathrm{Var}\{\Delta(\hat\eta)\}$,
\[
\mathrm{Var}(\hat\tau) = \frac{\E\,\mathrm{Var}(\psi_{\hat\eta})}{n_E} + \sigma_\Delta^2 + R_{n_E}, \qquad |R_{n_E}| \le \frac{2\sigma_\Delta C_1}{n_E} + \frac{C_1^2}{n_E^2} + \frac{C_2}{n_E^{3/2}}.
\]
Proof in Appendix~\ref{app:dmlproofs}.
\end{theorem}

The sandwich estimates the first term and misses the variance of bias across fits\index{standard error!what it leaves out|textbf}: one fit supplies one draw of $\Delta(\hat\eta)$, from which its variance cannot be estimated. This conditions on one fit; it gives no lower bound over all estimators and leaves open a bootstrap that refits the nuisance. The second term is negligible beside the first exactly when $\sqrt{n_E}\,\sigma_\Delta \to 0$, the rate condition for fluctuations rather than means. Under cross fitting, conditioning on all $K$ fits conditions on a statistic of every observation, because each fold trains the other fits. Only an identity survives.

\begin{proposition}[Cross fitting: the identity that survives]\label{crit:prop:cfvar}
Let $\hat\tau$ be the pooled cross fitted estimator\index{cross fitting!variance from the random split} with folds $I_1, \dots, I_K$ of sizes $n_k$ and fits $\hat\eta^{(-k)}$ that omit $I_k$. Set $\bar\Delta_n = \sum_k (n_k/n)\,b(\hat\eta^{(-k)})\big/\sum_k (n_k/n)\{J + \|\delta_m^{(k)}\|^2\}$, and $U_n = \hat\tau - \tau - \bar\Delta_n$. If $\hat\tau$ and $\bar\Delta_n$ are square integrable (the second holds when $\|\delta_m\| + \|\delta_\ell\| \le 1$ for every fit, since then $|\bar\Delta_n| \le \max(1, |\tau|)/J$) then\index{cross fitting!variance from the random split}
\begin{equation}
\mathrm{Var}(\hat\tau) = \mathrm{Var}(\bar\Delta_n) + 2\,\mathrm{Cov}(\bar\Delta_n, U_n) + \mathrm{Var}(U_n).
\label{eq:crit:cfvar}
\end{equation}
\end{proposition}
\begin{proof}
Write $\hat\tau - \tau = \bar\Delta_n + U_n$ and expand the variance of the sum.
\end{proof}

The sandwich is meant to estimate $\mathrm{Var}(U_n)$ by analogy with Theorem~\ref{crit:thm:single}; the identity does not establish that it does. If it does, it omits $\mathrm{Var}(\bar\Delta_n) + 2\,\mathrm{Cov}(\bar\Delta_n, U_n)$, whose cross term has no fixed sign. The paper's sufficient condition for a negligible cross term bounds how far one observation moves the fits. For the lasso in the worst case it requires $s\log p = o(\sqrt n)$; at $n = 300$, $p = 400$, $s = 5$, the two sides are about 30 and 17. Simulations know $\tau$ and both nuisances, so all four variance terms can be computed.

Take the simplest design in which they are computable, with no sparsity and no selection: $X \sim N(0, I_{20})$ with 20 irrelevant controls, treatment $T = \nu$ and outcome $Y = \tau T + \varepsilon$ with standard normal $\nu$ and $\varepsilon$ and $\tau = 1$, least squares nuisance fits, $n = 300$ and 30,000 replications. The cross term, as a multiple of $\mathrm{Var}(\bar\Delta_n)$, is $+0.02$ for a single split, which is zero within Monte Carlo error and consistent with the bound of Theorem~\ref{crit:thm:single}. Under cross fitting it is $+1.42$ at $K = 2$, $+0.21$ at $K = 3$, $-1.10$ at $K = 5$ and $-1.74$ at $K = 10$ (Exercise~\ref{ex:chcrit:7} checks the identity at $K = 5$). At two folds the omitted variance is two and a half times $\mathrm{Var}(\bar\Delta_n)$, and at five folds it is cancelled and reversed. Cross fitting induces a negative covariance between the realized bias and the noise around it, and $K$ is the dial. For a cross validated lasso at $p = 400$, with 2,000 replications a cell, the cross term at $K = 5$ is $-0.45$, $-0.39$, $-0.38$ and $-0.95$ times $\mathrm{Var}(\bar\Delta_n)$ at $s = 5$, 15, 30 and under weak overlap, four to eleven Monte Carlo standard errors from zero. At $K = 2$ and at a single split it is small and of either sign. In the dense design, where nothing covers, it stays within about two standard errors of zero at every fold count.

The omitted variance depends on $K$, while the reported error barely responds. Adding it closes only part of the gap: at $s = 30$ and five folds in the 2,000 replication run, adding $\mathrm{Var}(\bar\Delta_n) + 2\,\mathrm{Cov}(\bar\Delta_n, U_n) = 1.1\times 10^{-3}$ raises the reported error from 0.052 to 0.061 against a Monte Carlo spread of 0.073. The sandwich also falls short of $\mathrm{Var}(U_n)$, the block it should estimate: $2.7\times 10^{-3}$ against $4.2\times 10^{-3}$. We have no account of that gap.

Two procedures try to integrate over the fit, and they behave differently. The repeated split median of Chernozhukov and coauthors (2018), with the variance inflated by the spread across twenty splits, covers better than a single cross fit in all four designs where anything covers (Table~\ref{tab:crit:cmp}), at a cost of 3 to 11 percent in width, and moves the ratio of reported error to spread from 0.91, 0.81, 0.70 and 0.91 to 0.99, 0.93, 0.88 and 1.04. It repairs the variance and not the centre: under weak overlap coverage rises only from 0.557 to 0.603, because the error there is bias and the median of biased estimates is biased. The multiplier bootstrap\index{multiplier bootstrap|textbf} resamples the score at one fixed cross fit without refitting the nuisance (Mammen weights, 999 draws, percentile interval) and changes almost nothing: it covers 0.3 to 0.7 points below the single fit it resamples (McNemar $p$ of $2\times 10^{-4}$, 0.10, $9\times 10^{-4}$ and 0.004; Exercise~\ref{ex:chcrit:8}). Resampling the score is not integrating over the fit.

A bootstrap that does integrate over the fit is available at a price, and it overshoots. In a run of our own on 1,000 data sets per design ($n = 300$, $p = 400$, $\tau = 1$), we resampled the observations with replacement 499 times, refitted both cross validated lassos inside a fresh five fold cross fit in every resample, with all copies of an observation in the same fold and in the same inner validation fold, recomputed the estimate, and took the 2.5 and 97.5 percent points of the 499 estimates as the interval. On these data sets the single cross fit covers 0.916, 0.832 and 0.756 at $s = 5$, 15 and 30, and the refitting bootstrap 0.989, 0.998 and 0.984, with intervals 42, 70 and 126 percent wider: the mean squared bootstrap standard deviation is 1.7, 1.9 and 2.8 times the variance of the estimate across data sets. It is conservative, not calibrated. Under weak overlap it covers 0.607 against 0.536 for the single fit and 0.605 for the repeated split median, because the error there is bias, and in the dense design every procedure misses in every replication, with biases of 0.80 to 0.81 against mean interval widths of 0.16 to 0.30. There the bootstrap also understates the spread, with a mean squared standard deviation of $6.1\times 10^{-3}$ against a variance of $8.5\times 10^{-3}$.

\section{What no learner can repair}\label{sec:crit:repair}

Take eight equal strata of households, with propensities spaced evenly from 0.06 to 0.50 and effects from 2.0 down to 0. Give the partially linear score the true nuisance functions, so that no fitting error exists, and run it on $n = 4{,}000$ households for 500 replications. It returns 0.7701. The average effect is 1.0000. The interval covers 0.7706 on 0.960 of the replications and the average effect on none, and the gap does not shrink with $n$. The augmented score at the same true nuisances returns 0.9972 and covers the average effect on 0.966. Chapter~\ref{ch:dml} works this design by hand (Exercise~\ref{ex:chcrit:10} repeats the arithmetic), and Proposition~\ref{dml:prop:wate} explains the number: the weight $e(1 - e)$ is largest at a propensity of one half, so the stratum with the biggest effect and the weakest overlap gets the smallest weight. No learner can change that.

The weighting is not a defect in itself. Chapter~\ref{ch:dml} gives the case for choosing it, including trimming (Crump and coauthors 2009), and when the outcome variance is the same in both arms and at every covariate value, overlap weighting, whose weight is exactly $e(1 - e)$, has the smallest asymptotic variance among weighted average effects (Li, Morgan and Zaslavsky 2018). The objection is that DML arrives at this estimand silently: an analyst who calls a DML routine for the average treatment effect receives a different population without being told.

\paragraph{The efficiency bound and overlap.}\label{sec:dml:bound} Hahn (1998) gave the semiparametric efficiency bound\index{efficiency bound|textbf}\index{semiparametric efficiency bound|see{efficiency bound}} for the average effect. Write $\sigma_t^2(x) = \mathrm{Var}\{Y(t) \mid X = x\}$, $\tau(x) = \mu_1(x) - \mu_0(x)$ and $\bar\tau = \E\,\tau(X)$. The bound is
\begin{equation}
V_{\mathrm{eff}} = \E\left\{\frac{\sigma_1^2(X)}{e(X)} + \frac{\sigma_0^2(X)}{1 - e(X)} + \bigl(\tau(X) - \bar\tau\bigr)^2\right\}.
\label{eq:dml:bound}
\end{equation}
The first two terms are the outcome noise in each arm, inflated by the inverse probability of landing there, and they are what weak overlap does to any estimator of the average effect. In the example of Figure~\ref{fig:plmweights}, with $X \sim N(0,1)$, $\tau(x) = 1 + x$ and $e(x) = \expit(1 + \gamma x)$, the expectation $\E\{1/(1 - e(X))\}$ rises from 3.7 at $\gamma = 0$ to 246 at $\gamma = 3$, so with unit noise variances the best achievable variance rises about 46 times. The average effect is 1 for every $\gamma$, while the partially linear estimand is 0.70 at $\gamma = 1$, thirty percent lower, because the units it weights most, with propensity near one half, have below average effects. Whether a posterior attains the bound is a question about a Bernstein von Mises theorem for the average effect, taken up in Section~\ref{sec:crit:bayes} and in Appendix~\ref{app:theory}.

\begin{figure}[ht]
\centering
\includegraphics[width=\linewidth]{plm_weights.pdf}
\caption{Left: the partially linear estimand and the average effect as selection strengthens, for $\tau(x) = 1 + x$ and $e(x) = \expit(1 + \gamma x)$. Right: the three components of the efficiency bound \eqref{eq:dml:bound} with unit noise variances, on a log scale.}
\label{fig:plmweights}
\end{figure}

The Robinson denominator estimates $J$. At $J = 0$, $Y = \tau m(X) + g(X) + \varepsilon$, so any change in $\tau$ is absorbed by $g$ and the effect is unidentified. No procedure can approach this boundary safely\index{Dufour impossibility result|textbf}.

\begin{proposition}[No uniformly valid interval of finite length]\label{crit:prop:dufour}
Fix $n$. Let $\mathcal P$ be a class of laws of the partially linear model that contains, for each $j \in (0, \bar j\,]$ and each $\tau \in \mathbb R$, a law $P_{j,\tau}$ with $\E\,\mathrm{Var}(T \mid X) = j$ and effect $\tau$, and suppose that for every $L < \infty$
\[
\lim_{j \downarrow 0}\ \sup_{|\tau|, |\tau'| \le L}\ \bigl\|P_{j,\tau}^{\otimes n} - P_{j,\tau'}^{\otimes n}\bigr\|_{\mathrm{TV}} = 0.
\]
Let $\mathcal K_n$ be a confidence set with $\inf_{P \in \mathcal P} \Prob_P(\tau(P) \in \mathcal K_n) \ge 1 - \alpha$ for some $\alpha < 1$, whose graph is jointly measurable (as it is for any interval with measurable endpoints). Then $\sup_{P \in \mathcal P}\E_P\,\mathrm{Leb}(\mathcal K_n) = \infty$\index{Dufour impossibility result}. The Gaussian family $T = m(X) + \nu$, $\nu \sim N(0, j)$, $Y = \tau\nu + h(X) + \varepsilon$, $\varepsilon \sim N(0, \sigma_\varepsilon^2)$ satisfies the condition, with $\|P_{j,\tau}^{\otimes n} - P_{j,\tau'}^{\otimes n}\|_{\mathrm{TV}} \le L\sqrt{nj}/\sigma_\varepsilon$ for $|\tau|, |\tau'| \le L$. Proof in Appendix~\ref{app:dmlproofs}.
\end{proposition}

The proof is a Fubini argument in the manner of Gleser and Hwang (1987) and Dufour (1997): the expected length of $\mathcal K_n$ is the integral over $\tau'$ of the probability that it covers $\tau'$, and near $j = 0$ that probability is close to $1 - \alpha$ over a window as wide as one likes. The proposition carries two safeguards.

First, it does not contradict Chernozhukov and coauthors (2018), who fix $c > 0$ and require $\E\,\mathrm{Var}(T \mid X) \ge c$ throughout their class, on which the sandwich interval is uniformly valid. Here $J$ is positive at every law and has no common lower bound. That bound is an assumption, and the data cannot certify it.

Second, a fixed design cannot exhibit uniformity failure. The vanishing overlap sweep shows loss of calibration before the boundary. In our sweep ($s = 10$, five folds) the treatment residual standard deviation takes the values 1.0, 0.6, 0.35, 0.2 and 0.1. The width of the DML interval with a cross validated lasso rises by a factor of 10.1 and its bias from 0.034 to 0.907, a factor of 27, so coverage falls from 0.848 to 0.610. The reported error is right, with a ratio to the Monte Carlo spread between 0.85 and 1.05 down the sweep, and the centre is wrong. The deterioration belongs to the cross validated lasso. Post double selection at the theory penalty on the same data sets has bias $0.006$, $-0.006$ and $-0.033$ at residual standard deviations 0.35, 0.20 and 0.10, against 0.208, 0.428 and 0.907 for DML, with coverage 0.956, 0.946 and 0.972 against 0.704, 0.626 and 0.610, and a horseshoe nuisance at five folds does not show the deterioration either. The theory penalty succeeds only where it recovers the support, which it does on 1,493 of 1,500 fits at $s = 10$ and on none at $s = 30$, and no data driven check separates the two cases.

The Bayesian analogue is not exempt: a credible set of finite expected length over the same class is not uniformly valid either. In the same sweep the reduced form posterior of Section~\ref{sec:crit:bayes} widens from 0.247 to 2.631 and its coverage rises from 0.806 to 0.970, while DML's falls. Both intervals widen by about the same factor, 10.1 for DML and 10.7 for the posterior. What separates them is the centre: the bias grows 27 times for DML and 1.9 times for the posterior, from 0.068 to 0.126.

\begin{remark}[No certificate]\label{crit:rem:cert}
The bias functional $b$ is a combination of an inner product and a squared norm of the unknown nuisance errors. Certifying the rate condition therefore means estimating those two functionals. In the regime $s \gg \sqrt n/\log p$ the minimax rate for functionals of this kind is of order $s\log p/n$ (Guo, Wang, Cai and Li 2019), which is the order of $b$ itself, and the impossibility of adaptive confidence intervals (Cai and Guo 2017) is the same phenomenon one level down. This is an argument and not a theorem, since no matching lower bound for the composite functional is proved. For smooth nuisance functions the barrier is known: root $n$ estimation of the average effect needs roughly $\beta_e + \beta_\mu \ge d/2$, for propensity and outcome smoothness $\beta_e$ and $\beta_\mu$ in dimension $d$, and below that no interval of width $O(n^{-1/2})$ is valid (Robins, Tchetgen Tchetgen, Li and van der Vaart 2009; Robins, Li, Mukherjee, Tchetgen Tchetgen and van der Vaart 2017; Balakrishnan, Kennedy and Wasserman 2023).
\end{remark}

No general test of the rate condition is known. Unconfoundedness and no interference cannot be tested from observational data. If unconfoundedness fails, the average effect lies in a set whose size depends on the failure (Manski 1990); DML maintains the assumption and reports a point, saying nothing about that set. A Bayesian prior states how strongly an unmeasured confounder might act and visibly carries the unidentified part (Chapter~\ref{ch:tianpearl}).

\section{The same anatomy in weighting}\label{sec:crit:weights}

Replace the partially linear score by the augmented score \eqref{eq:aipw} of Chapter~\ref{ch:dml} and the estimand problem goes away. The error problem does not. With $\hat e$ and $\hat\mu_t$ the fitted propensity and arm regressions, the mean of the score at the fitted nuisances, derived there, is a sum of two products of the propensity error with the outcome error in each arm.
Each term is bounded by a Cauchy and Schwarz product, and the alignment question returns, since a covariate that raises the propensity usually raises the outcome too. Polson and Sokolov (2026b) state that the sharp rate condition transfers to this score with a weighted cosine between the propensity error and the outcome errors and a factor $1/\hat e_{\min}$ in the bound, where $\hat e_{\min}$ is the smallest value of $\hat e(x)$ and $1 - \hat e(x)$, so that weak overlap multiplies the alignment penalty instead of acting separately from it. The proof is in their supplement and is not reproduced here.

The augmented score is doubly robust\index{double robustness}: it is consistent if either the propensity or the outcome model is correct (Robins, Rotnitzky and Zhao 1994). Three qualifications apply. Both models can be wrong: in high dimensions both are approximations and the guarantee is silent, and when both are moderately wrong the estimator can be worse than regression adjustment alone, because the weighting term amplifies the outcome model's errors where the estimated propensity is small (Kang and Schafer 2007; Chapter~\ref{ch:ipw}). The guarantee is asymptotic, and in finite samples, especially with propensities near zero or one, the remainder terms can dominate. And a posterior does not become doubly robust by being Bayesian. It can put a direction for the propensity score into its prior, as the Bayesian causal forest and the coupled horseshoe of Chapter~\ref{ch:horseshoe} do, which lowers its sensitivity to confounding in that direction. Bayesian AIPW, in the same chapter, combines a posterior with the augmented score explicitly and inherits double robustness from the score.

\section{Where Bayes helps and where it does not}\label{sec:crit:bayes}

Suppose the analyst, persuaded by the last four sections, writes the Bayesian analysis a well trained Bayesian would write: $Y = \tau T + X^\top\beta + \varepsilon$, a horseshoe prior on $\beta$ and a flat prior on $\tau$. At $s = 30$ confounders ($n = 300$, $p = 400$) its posterior mean averages 1.96 for a true effect of 1, its 95 percent interval is 0.073 wide, and the interval covers the truth in none of 5,000 replications\index{credible interval!frequentist coverage of|textbf}\index{coverage!frequentist, of credible intervals|textbf}. The DML arms report intervals about 0.2 wide. The failure has a name and a limit.

\begin{proposition}[Complete shrinkage bias]\label{crit:prop:ric}\label{crit:prop:default}
Let $T = X^\top\gamma + \nu$ and $Y = \tau T + X^\top\beta + \varepsilon$, with $\E X = 0$, $\mathrm{Var}(X) = \Sigma$, $\E \nu = 0$, $\mathrm{Var}(\nu) = \sigma_\nu^2$, $\nu$ uncorrelated with $X$ and $\E(\varepsilon \mid T, X) = 0$ with $\E \varepsilon^2 < \infty$. Shrinking $\beta$ entirely to zero\index{default prior} and regressing $Y$ on $T$ alone converges in probability, under independent sampling as $n \to \infty$, to
\[
\tau_\infty = \tau + \frac{\gamma^\top\Sigma\beta}{\gamma^\top\Sigma\gamma + \sigma_\nu^2}.
\]
If $\Sigma = I_p$ and $\gamma_j = \beta_j = c \ne 0$ on a set of size $s$ and zero elsewhere, the bias is $sc^2/(sc^2 + \sigma_\nu^2)$, which increases in $s$ and tends to one.
\end{proposition}
\begin{proof}
Since $\E(\varepsilon T) = 0$, the univariate coefficient converges to
\[
\frac{\mathrm{Cov}(T,Y)}{\mathrm{Var}(T)} = \tau + \frac{\gamma^\top\Sigma\beta}{\gamma^\top\Sigma\gamma + \sigma_\nu^2}.
\]
With $\Sigma = I_p$ and the stated coefficients, both quadratic forms equal $sc^2$.
\end{proof}

The limit is omitted variable bias\index{omitted variable bias}, the bias when shrinkage is complete. In the symmetric design, dropping $k$ of the $s$ confounders and keeping the rest gives the smaller bias $kc^2/(kc^2 + \sigma_\nu^2)$, so there no selection does worse than dropping them all. In other designs the limit is not a ceiling: with $\Sigma = I_2$, $\gamma = (1, 1)^\top$, $\beta = (1, -1)^\top$ and $\sigma_\nu = 1$, complete shrinkage has no bias, while dropping the first confounder alone leaves a bias of $1/2$. As the prior pulls $\beta$ toward zero, the treatment coefficient, the only unpenalized parameter, absorbs the confounding that the shrunken coefficients no longer explain, which is regularization induced confounding\index{regularization induced confounding} (Hahn, Carvalho, Puelz and He 2018). The sampling variance of the misspecified regression is of order $1/(ns)$ in this design, so the interval narrows as the bias grows. At $c = \sigma_\nu = 1$ the limiting bias is 0.83, 0.94 and 0.97 at $s = 5$, 15 and 30. The naive Bayes biases in our designs are 0.032, 0.362 and 0.956, about 4, 39 and 99 percent of the limit, and the interval widths are 0.224, 0.490 and 0.073 (at $s = 15$ the posterior has two modes, near 1 and near 2, and the pooled interval spans both when the chains of a replication disagree): the posterior is most confident where it is most wrong. In this design the limit is a ceiling and not a prediction for finite $n$. The harm comes from applying a prior chosen for its estimation properties in coordinates that ignore the design, and the reduced form below applies the same family of prior to different coefficients.

The remedy is a change of coordinates. Following DiTraglia and Liu (2025), put horseshoe priors on the coefficients of the two nuisance regressions, $m(x) = x^\top\gamma$ and $\ell(x) = x^\top\varsigma$ with $\varsigma = \tau\gamma + \beta$, instead of on the structural coefficient $\beta$. Chapter~\ref{ch:horseshoe} takes the other route and couples the two regressions in the structural coordinates. Table~\ref{tab:crit:head} gives three of our five designs.

The reduced form model is $T = X\gamma + \nu$ and $Y = X\varsigma + \tau\nu + \varepsilon$, with $\nu \sim N(0, \sigma_\nu^2 I)$ and $\varepsilon \sim N(0, \sigma_\varepsilon^2 I)$ independent and $\varsigma = \tau\gamma + \beta$. The prior is a horseshoe on $\gamma$ with scales proportional to $\sigma_\nu$ and a horseshoe on $\varsigma$ with scales proportional to $\sigma_\varepsilon$, each with its own half Cauchy local and global scales, a flat prior on $\tau$ and Jeffreys priors on the two variances, and the sampler is Gibbs. Given a nuisance draw $\eta = (\gamma, \varsigma, \sigma_\nu, \sigma_\varepsilon)$, let $\nu_\gamma = T - X\gamma$ and $r_\varsigma = Y - X\varsigma$. The effect is then normal with mean $a_\eta = \nu_\gamma^\top r_\varsigma/\nu_\gamma^\top\nu_\gamma$ and variance $V_\eta = \sigma_\varepsilon^2/\nu_\gamma^\top\nu_\gamma$, a Robinson regression at the drawn nuisance. By the law of total variance, $\mathrm{Var}(\tau \mid D) = \E(V_\eta \mid D) + \mathrm{Var}(a_\eta \mid D)$. The first term plays the part of the sandwich, and the second is the spread of the centre across nuisance draws, the posterior counterpart of $\sigma_\Delta^2$ in Theorem~\ref{crit:thm:single}. The two are not shown to be equal, since one is a posterior spread and the other a frequentist variance.

\begin{table}[ht]
\centering
\small
\caption{Head to head in three of our five designs. $n = 300$, $p = 400$, $\tau = 1$, nominal 95 percent intervals, seeds shared across arms; 5,000 replications a cell in the sparse and weak overlap blocks (the data sets of Table~\ref{tab:dml:sims}) and in the dense block 4,406 for the lasso row, 5,000 for the horseshoe rows and 1,000 for naive Bayes and the posterior. In the sparse and weak overlap blocks the reduced form posterior and naive Bayes run four chains of 10,000 sweeps per replication, doubled until split $\hat R$ for $\tau$ is at most 1.01 and its effective sample size at least 400, up to 80,000 sweeps a chain; every reduced form replication passed, naive Bayes at $s = 15$ in 59 percent, because its posterior there has two modes. The horseshoe nuisance is the posterior mean from 1,200 sweeps. DML rows carry their fold count; the reduced form posterior is the default Bayes arm. The last two columns divide the mean reported standard error by the Monte Carlo standard deviation of the estimator and by its root mean squared error. The dense posterior row uses four chains per replication, run until split $\hat R$ for $\tau$ is at most 1.01 and its effective sample size at least 400, which took up to 800,000 sweeps a chain; 103 of the 1,000 replications had not passed by then and 7 that passed report a standard error above 1, and the row is over the other 890.}
\label{tab:crit:head}
\begin{tabular}{@{}lrrrr@{}}
\toprule
 & bias & coverage & s.e./s.d. & s.e./RMSE \\
\midrule
\multicolumn{5}{@{}l}{\textit{sparse, $s = 15$}} \\
DML, lasso, $K = 5$ & 0.037 & 0.826 & 0.81 & 0.70 \\
DML, horseshoe, $K = 5$ & $-0.012$ & 0.895 & 0.85 & 0.83 \\
DML, horseshoe, $K = 2$ & $-0.069$ & 0.587 & 0.49 & 0.43 \\
naive Bayes & 0.362 & 0.446 & 0.56 & 0.36 \\
reduced form posterior & 0.076 & 0.787 & 1.02 & 0.66 \\
\addlinespace
\multicolumn{5}{@{}l}{\textit{weak overlap, $s = 15$, $\sigma_\nu = 0.3$}} \\
DML, lasso, $K = 5$ & 0.317 & 0.557 & 0.91 & 0.47 \\
DML, horseshoe, $K = 5$ & $-0.024$ & 0.908 & 0.87 & 0.86 \\
DML, horseshoe, $K = 2$ & $-0.040$ & 0.841 & 0.74 & 0.73 \\
naive Bayes & 0.987 & 0.000 & 0.48 & 0.02 \\
reduced form posterior & 0.148 & 0.906 & 1.01 & 0.85 \\
\addlinespace
\multicolumn{5}{@{}l}{\textit{dense confounding}} \\
DML, lasso, $K = 5$ & 0.799 & 0.000 & 0.46 & 0.05 \\
DML, horseshoe, $K = 5$ & 0.677 & 0.000 & 0.56 & 0.09 \\
DML, horseshoe, $K = 2$ & 0.783 & 0.000 & 0.54 & 0.06 \\
naive Bayes & 0.935 & 0.000 & 0.99 & 0.02 \\
reduced form posterior & 0.938 & 0.000 & 1.30 & 0.03 \\
\bottomrule
\end{tabular}
\end{table}

The ratio of reported error to spread is the column to read first, because it is the one Theorem~\ref{crit:thm:single} is about. The reduced form posterior reports 1.05, 1.02, 1.00 and 1.01 of its own spread at $s = 5$, 15, 30 and under weak overlap, with converged chains in every one of the 5,000 replications. DML with the cross validated lasso at five folds reports 0.91, 0.81, 0.70, 0.91 and 0.46 at $s = 5$, 15, 30, under weak overlap and in the dense design, and with a horseshoe nuisance 0.95, 0.85, 0.73, 0.87 and 0.56, falling to between 0.45 and 0.74 at two folds. This is the one place where the evidence is unambiguous: a posterior over the nuisance pair integrates over the fit by construction, and its reported error does not omit the dispersion term. It is no certificate. In the dense design, where every coordinate is nonzero and no learner makes the nuisance error vanish, the posterior reports 1.30 times its spread, is biased by 0.94 and covers in none of the 890 replications of the table. In 103 of the 1,000 its four chains, started from widely dispersed values, had not converged after 800,000 sweeps, and their intervals are so wide (median width 175) that 71 of them cover; 7 more converged with standard errors from 1.25 to 610, and the table leaves all 110 out. The ratio that governs coverage is the last column, reported error over root mean squared error, and it is 0.03 there.

On the estimate itself the picture is mixed, and DML wins the sparse designs. Its bias at $s = 5$, 15 and 30 is 0.017, 0.037 and 0.050 against 0.041, 0.076 and 0.106 for the posterior, paired differences of 0.025, 0.039 and 0.056, each more than 60 standard errors from zero. Coverage separates them at $s = 15$ and 30 and not at $s = 5$: 0.917, 0.826 and 0.734 against 0.913, 0.787 and 0.704 (McNemar $p$ of 0.25, $4\times 10^{-13}$ and $6\times 10^{-6}$), and both fall well below nominal by $s = 30$. Under weak overlap\index{overlap!weak} the order reverses: the posterior covers on 0.906 of the replications against 0.557, on 1,776 discordant pairs to 30. DML with a horseshoe nuisance at five folds is the best centred arm, with absolute bias 0.011, 0.012, 0.021 and 0.024 in the four designs where anything covers. Its coverage, 0.928, 0.895, 0.827 and 0.908, is above the lasso arm's in all four designs (McNemar $p = 0.008$ at $s = 5$ and below $10^{-27}$ in the other three) and above the posterior's by 1.5, 10.8 and 12.4 points at $s = 5$, 15 and 30 ($p = 0.001$ and below $10^{-45}$), and under weak overlap it matches the posterior, 0.908 against 0.906 ($p = 0.74$). Cutting its folds to two cuts its coverage to 0.577, 0.587, 0.671 and 0.841. In the dense design, where nothing covers (Table~\ref{tab:crit:head}), no interval should be reported.

To separate DML from the cross validated lasso inside it, Table~\ref{tab:crit:cmp} runs four more procedures on the same data sets by seed replay.

\begin{table}[ht]
\centering
\small
\caption{Four more procedures on the data sets of Table~\ref{tab:crit:head} ($n = 300$, $p = 400$, $\tau = 1$; the dense design is omitted), replayed on the same seeds. The first row is DML with the cross validated lasso at five folds. Post double selection refits least squares on the union of the variables chosen by lasso regressions of the outcome and of the treatment on the controls, with the lasso penalty set by theory or by cross validation. The repeated split median and the multiplier bootstrap are the two procedures of Section~\ref{sec:crit:se}. Nominal 95 percent intervals, 5,000 replications a cell.}
\label{tab:crit:cmp}
\begin{tabular}{@{}lrrr@{}}
\toprule
 & bias & coverage & mean width \\
\midrule
\multicolumn{4}{@{}l}{\textit{sparse, $s = 5$}} \\
DML, cross validated lasso & 0.017 & 0.917 & 0.209 \\
post double selection, theory penalty & $-0.002$ & 0.950 & 0.216 \\
post double selection, CV penalty & $-0.086$ & 0.669 & 0.249 \\
repeated split median & 0.017 & 0.936 & 0.215 \\
multiplier bootstrap & 0.017 & 0.911 & 0.208 \\
\addlinespace
\multicolumn{4}{@{}l}{\textit{sparse, $s = 15$}} \\
DML, cross validated lasso & 0.037 & 0.826 & 0.205 \\
post double selection, theory penalty & 0.798 & 0.010 & 0.127 \\
post double selection, CV penalty & $-0.110$ & 0.629 & 0.285 \\
repeated split median & 0.038 & 0.870 & 0.218 \\
multiplier bootstrap & 0.037 & 0.823 & 0.204 \\
\addlinespace
\multicolumn{4}{@{}l}{\textit{sparse, $s = 30$}} \\
DML, cross validated lasso & 0.050 & 0.734 & 0.201 \\
post double selection, theory penalty & 0.960 & 0.000 & 0.066 \\
post double selection, CV penalty & $-0.120$ & 0.651 & 0.337 \\
repeated split median & 0.051 & 0.820 & 0.224 \\
multiplier bootstrap & 0.050 & 0.728 & 0.200 \\
\addlinespace
\multicolumn{4}{@{}l}{\textit{weak overlap, $s = 15$, $\sigma_\nu = 0.3$}} \\
DML, cross validated lasso & 0.317 & 0.557 & 0.687 \\
post double selection, theory penalty & 0.924 & 0.034 & 0.138 \\
post double selection, CV penalty & $-0.134$ & 0.867 & 0.925 \\
repeated split median & 0.316 & 0.603 & 0.723 \\
multiplier bootstrap & 0.317 & 0.550 & 0.684 \\
\bottomrule
\end{tabular}
\end{table}

Post double selection\index{double selection|textbf} (Belloni, Chernozhukov and Hansen 2014) at the theory penalty is the best procedure in the table at $s = 5$, with bias $-0.002$ and coverage 0.950 against DML's 0.917 (McNemar $p = 2\times 10^{-20}$). At $s = 15$ and $s = 30$ it fails completely, with biases of 0.798 and 0.960, because the theory level is calibrated to a sparsity the design does not have and the joint lasso keeps 8.9 and 2.2 coordinates. At the cross validated penalty it keeps 37.8, 89.4 and 146.3 coordinates and covers below DML in all three sparse designs. Which tuning is good depends on the sparsity, which the data cannot certify. Under weak overlap post double selection at the cross validated penalty covers on 0.867 against DML's 0.557 and below the posterior's 0.906 (McNemar $p = 5\times 10^{-10}$), while DML with a horseshoe nuisance matches the posterior, so the weak overlap result is a statement about the cross validated lasso inside the score more than about Bayes against frequentism. The reduced form posterior's lowest coverage in these four designs is 0.704, against 0.557 for DML, 0.000 for post double selection at the theory penalty, 0.629 at the cross validated penalty, 0.603 for the repeated split median\index{cross fitting!repeated splits} and 0.550 for the multiplier bootstrap\index{multiplier bootstrap}. Post double selection at the theory penalty is ahead of the posterior at $s = 5$ and far behind it at the other three designs, while DML and the repeated split median cover at least as well as the posterior at $s = 5$, 15 and 30 and worse under weak overlap. Which case applies turns on sparsity and overlap, which the data do not reveal, so the argument is about worst cases and not about averages. Only DML with a horseshoe nuisance at five folds, whose lowest coverage is 0.827, does better in the worst case, and it costs 19.1 seconds a data set.

Cross fitting the reduced form posterior makes it worse. Coverage in the four designs where anything covers falls from 0.913, 0.787, 0.704 and 0.906 for the pooled posterior on 5,000 replications to 0.737, 0.351, 0.066 and 0.754 at five folds, on 4,976, 4,991, 4,999 and 5,000 replications. Each fold's posterior is sampled by four chains, the held out estimate is combined across the folds at every draw, and the four chain rule of Table~\ref{tab:crit:head} is applied to that combined estimate; it passes in 88, 98, 98 and 99.5 percent of the replications, so the first figure is the least secure. At two folds the sampler mixes far more slowly. With four chains per fold from overdispersed starts, $\tau$ drawn jointly with $\varsigma$, and 100,000 sweeps doubled up to 800,000, on 500 replications per design, coverage is 0.55, 0.04, 0.01 and 0.33, and the chains pass the same rule in 466, 292, 494 and 307 of the 500. At $s = 15$ and under weak overlap about two replications in five have not converged, so those two coverages pool converged and unconverged chains, and any of the four carries a Monte Carlo standard error of up to 0.022. The direction is not in question: the cross fitted posterior is biased downward in every design, by 0.08 to 0.24 at five folds and by 0.12 to 0.45 at two, and in the replications whose chains converged as well. The pooled posterior's small bias is a cancellation between a negative pure bias term and a positive dependence term, and there is no argument for why the two should track. Its bias is positive in every design, so the mechanism of Proposition~\ref{crit:prop:ric} survives the change of coordinates in weakened form. Its centring in these designs is an empirical finding with no theorem behind it.

The orthogonal score is the efficient score for $\tau$ in the partially linear model when the outcome noise is homoskedastic (otherwise efficiency needs variance weights), and the augmented score is the efficient influence function for the average effect. A Bayesian model can share their first order behaviour when its prior contains the direction those scores use, the propensity score\index{propensity score}. Ray and van der Vaart (2020) prove a semiparametric Bernstein von Mises theorem\index{Bernstein von Mises theorem} for the mean of a binary outcome missing at random, one arm of an average effect, under smoothness conditions that weaken when the Gaussian process prior includes a component in that direction (Appendix~\ref{app:theory} states the scope), and the coupled horseshoe of Chapter~\ref{ch:horseshoe} and the Bayesian causal forest of Hahn, Murray and Carvalho (2020) put the propensity into the prior for the same reason. We prove nothing of the kind for the horseshoe, and the measurements above are the evidence we have: modelling the treatment equation together with the outcome, as the reduced form does, shrinks the regularization induced confounding of Proposition~\ref{crit:prop:ric} without ending it. Orthogonality is a property a Bayesian model can aim at by construction, and a posterior adds the rest of the distribution.

Cheapness is the one advantage we grant DML without qualification. In a one core timing run of our own ($n = 300$, $p = 400$, ten data sets of the $s = 15$ design), a cross fitted lasso took 1.0 seconds on average against 9.8 for one chain of 2,000 sweeps of the reduced form sampler, and it needs no convergence diagnostics, though the analyst still chooses the score, the learners, their tuning and the folds. The cheapness belongs to the lasso and not to the template: DML with a horseshoe nuisance at five folds took 19.1 seconds. The posterior of Table~\ref{tab:crit:head} runs four chains and extends them until they pass the convergence checks, which took 194 to 227 seconds, about 200 times the lasso; naive Bayes took 96 seconds on seven of the data sets and about 1,440 on the other three.

\paragraph{What four people get.}\label{sec:crit:applied} The case for a posterior rests less on the point estimate of an average effect than on what a decision needs. Under the model, and with a prior the analyst can defend, one fit supplies each of the following; a DML analysis would need a separate score for each, and none of them comes out as a probability. A regulator gets the probability that the effect exceeds a clinically meaningful threshold, the effect in each prespecified subgroup and a predictive distribution for the next patient. A ministry with a fixed budget needs the conditional effect for each household to rank them, its uncertainty to avoid chasing noise, and the distribution of gains to judge equity. In the 401(k) data of Section~\ref{sec:deephd} one quantile network, refitted 200 times under Bayesian bootstrap weights, gives the effect in every income quintile, from about three thousand dollars to eighteen, and the whole conditional distribution of assets behind each. DML with the augmented score also gives quintile effects from one set of nuisance fits (Chapter~\ref{ch:k401}), but not the distributions. A court asks whether this plaintiff's cancer was caused by smoking, a probability of causation that data bound but do not identify; the answer is a posterior inside the bounds with its prior stated, and DML has no object here, since no average identifies a probability that concerns one person (Chapter~\ref{ch:tianpearl}). An economist modelling a merger needs the price effect as a function of demand elasticities, conduct and costs, and then the probability that prices rise by more than a competition authority's threshold (Chapter~\ref{ch:blp}); a point estimate with a standard error is where that calculation starts. The analyst pays with a model and a prior. The predictive check of Chapter~\ref{ch:k401} shows replicated 401(k) wealth from a Gaussian outcome model to be far too symmetric, while a Student $t$ error model with three degrees of freedom reproduces the bulk of the data much better. The prior must put the design in, or regularization induced confounding returns. DML has assumptions of its own, and none of them appears in its output.

\section{Three real data sets}\label{sec:crit:real}



\paragraph{401(k) eligibility.} The 1991 SIPP extract of Chernozhukov and Hansen (2004) has $n = 9{,}915$, eligibility as the treatment, net financial assets as the outcome and, in our basis, $p = 99$ controls (not the basis of Chapter~\ref{ch:k401}). Nine procedures that disagree sharply in the simulations, all built on the Gaussian partially linear score, return between 9.41 and 10.67 thousand dollars, a spread of 1.26, which is 0.94 of one least squares standard error. This is the regime DML was built for: $p/n$ is 0.01, the treated fraction is 0.37, and the mean squared treatment residual is 0.196, nowhere near zero. In this regime DML is a good choice and this chapter does not argue otherwise. The agreement is within a family, though. Three procedures that change the score or the likelihood give 5.81 (a Student $t$ outcome likelihood with four degrees of freedom, whose interval is 1.47 wide against 4.67 for the Gaussian one, which returns 9.56), 8.34 (the augmented score, standard error 1.31) and 7.16 (a stratified horseshoe over deciles of an estimated propensity, standard error 1.23). The $t$ likelihood estimates a robust centre and not a mean, and net financial assets have skewness 10.6 and excess kurtosis 184, so these three do not estimate the same functional as the nine. The three sit 1.1 to 3.6 thousand dollars below the lowest of the nine, which span 1.26 among themselves.

\paragraph{LaLonde.} LaLonde (1986) asked whether observational estimators recover the experimental effect of the National Supported Work demonstration when the randomized controls are replaced by a survey comparison group. We use the 185 treated units of the Dehejia and Wahba (1999) sample against the 2,490 units of the PSID comparison group, with ten covariates expanded to $p = 56$ terms, and three truths: the experimental effect on the treated, 1.79 thousand dollars, and zero for each of two placebos (the 260 randomized controls in place of the treated, and 1975 earnings as the outcome with the 1975 covariates withheld). Every procedure built on the partially linear score returns between $-0.23$ and $0.46$ against 1.79. We read this as the overlap weighted estimand of Proposition~\ref{dml:prop:wate}, since most treated units have propensities near one, where $\mathrm{Var}(T \mid X)$ is smallest. On the real data that reading cannot be separated from confounding. Every procedure, Bayesian or not, fails the placebo by between 2.4 and 7.3 thousand dollars, which suggests that the PSID comparison is not ignorable given the available covariates (Smith and Todd 2005 revisit the sensitivity of such comparisons), so the gap mixes the estimand with the failure of identification. Only a semi synthetic version of the design separates them. It keeps the real covariates of the 2,675 units and simulates the rest in 100 replications: treatment is Bernoulli with the propensity $e$ of a ridge logistic regression of the real treatment on the expanded basis, redrawn every replication; the untreated outcome is a gradient boosting\index{gradient boosting|textbf} fit (depth 3, 300 trees) of real 1978 earnings on the ten main covariates among the 2,490 comparison units, plus noise resampled from those units' residuals; and the effect is $1.794 + 1.5\{e(x) - \bar e_T\}$ with $\bar e_T = \E e^2/\E e$, so that the expected effect on the treated is the experimental benchmark and the effect rises with the propensity. There the partially linear DML rows cover the overlap weighted effect they target (1.31 in that design) on 0.94 of replications and the effect on the treated on 0.87. On real data the structural assumption fails before the estimator does.

\paragraph{IHDP.} Hill (2011) removed the treated children of nonwhite mothers from the Infant Health and Development Program\index{Infant Health and Development Program}, which leaves 747 children of whom 139 are treated and makes overlap incomplete by construction, kept the 25 real covariates and the real treatment, and simulated outcomes from Hill's linear surface: in each replication the untreated outcome is $b_0 + X^\top b$ plus standard normal noise, with $b_0$ drawn uniformly from $\{0, 1, 2, 3, 4\}$ and each of the 25 slopes drawn from the same five values with probabilities 0.5, 0.125, 0.125, 0.125 and 0.125, and the treated outcome adds a constant effect of 4 and its own standard normal noise. The target is the effect on the treated, 4 in every replication, and the setting is 100 replications. The partially linear DML score covers in none of them. The effect is constant, so every weighted average of it is 4 and Proposition~\ref{dml:prop:wate} is not the explanation. The mechanism of Theorem~\ref{crit:thm:atten} fits what follows, although its hypotheses are not checked here: the treatment is binary, so its residual variance is not constant in $X$, and nothing shows the cross channel to be negligible. In a separate run of 12 replications at effects of 0, 4 and 8, the estimate at 4 and 8 is 0.78 and 0.79 of the truth, and with an effect of zero there is nothing to attenuate. In another, of 20 replications, replacing the outcome learner by a post lasso refit, or by least squares on any of three bases, leaves the bias between $-0.88$ and $-1.12$. In the 12 replication run, replacing the cross validated $\ell_1$ logistic propensity on a 160 column basis by an unpenalized logistic on the 25 main effects removes almost all of it and leaves a ratio of 0.98. Cross validation shrank the propensity toward its mean, the error it left entered the treatment residual as measurement error, and the estimator returned the reliability ratio\index{reliability ratio} times the effect. Its own denominator reports the ratio: $\hat J_n$ is 0.180 under the penalized propensity and 0.147 under the unpenalized one, so if the second is close to $J$ the representation predicts $0.147/0.180 = 0.82$ against the 0.78 observed. This reads as the own square channel binding with the outcome regression essentially exact, and it is invisible from inside the procedure: the reported standard error is 1.07 times the estimator's spread, and the interval is centred at 3.12.

The Bayesian side has one resolved gain and one resolved loss. In the semi synthetic LaLonde design two posteriors for the effect on the treated, one from an outcome model with the propensity as a covariate and one from a prior that is exchangeable within propensity strata, cover at 0.97 and 0.99 against 0.88 for the doubly robust DML interval (exact McNemar $p = 0.022$ and $0.003$), on intervals a third wider and with no resolved difference in root mean squared error. On IHDP the procedures that target the effect on the treated are all accurate and 100 replications do not separate them, with one exception that runs against default Bayes, the reduced form posterior with the binary treatment treated as Gaussian: on the linear surface it covers at 0.90 and the doubly robust DML interval at 0.97 (exact McNemar $p = 0.039$).

\section{What this chapter does not establish}\label{sec:crit:limits}

The critique is theoretical and the evidence is simulated, and the gap between the two is where a reader should press. The designs are linear and Gaussian, with $n = 300$ and $p = 400$. The paper runs a cross validated lasso, a theory penalty lasso in one experiment and a horseshoe posterior mean inside DML; our forest, boosting and stack comparison uses one law that favours linear learners. Neural networks, propensity calibration, BART on IHDP and a binary treatment model inside the reduced form sampler are untried. The 5,000 replication comparisons have coverage Monte Carlo standard errors of at most 0.7 points; the paper's 500 replication runs have up to 2.2, so differences of two or three points are read only where the paired McNemar test resolves them. Cross validation never meets the support recovery of Proposition~\ref{crit:prop:alignment}, nothing delivers the negligible cross channel of Theorem~\ref{crit:thm:atten}, and the sandwich's shortfall against $\mathrm{Var}(U_n)$ remains unexplained (Section~\ref{sec:crit:se}). The reduced form sampler needs long chains: 117 of 20,000 sparse and weak overlap replications needed more than 10,000 sweeps a chain, one needed 40,000, and the dense and two fold failures reported in Section~\ref{sec:crit:bayes} qualify those coverages. Naive Bayes at $s = 15$ pools chains in different modes in 41 percent of replications after 80,000 sweeps. The horseshoe nuisance runs 1,200 sweeps with no convergence check. Rerun at 5,000 sweeps on the first 500 data sets of each design, its coverage moves by at most 0.042, except at $s = 5$ and two folds, where it rises from 0.59 to 0.74, so that entry depends on the length of the chain. Of the three real data sets two have a truth to check against, and on LaLonde every procedure fails the placebo. Four problems are open: a horseshoe result for a general design; a proof or a refutation of the shared false positive mechanism behind the alignment; a feasible frequentist estimate of the distribution of the bias functional, beyond the conservative refitting bootstrap of Section~\ref{sec:crit:se}; and a default prior for the dense regime, where the argument of Remark~\ref{crit:rem:cert} suggests that no interval may be available at any price.

\section{In practice}\label{sec:crit:practice}

The target, the checklist and the four conditions under which DML with the partially linear score is the right tool are in the In practice section of Chapter~\ref{ch:dml}. The evidence of this chapter adds the following.

\paragraph{What the analyst does differently.} Fit the nuisance equations at both the cross validated and the theory penalty and compare the selected supports; a large gap is a reason for caution, though this heuristic's operating characteristics are untested. With the orthogonal score, use repeated splits with the median and the variance inflated by their spread, and report the number of splits and the number of folds. Report the ratio of the reported standard error to any available spread and run the posterior as well; the two disagreeing is information, and the two agreeing is not a certificate. Report the effect stratum by stratum with the number of comparison units behind each, and run a placebo test where one exists. On the LaLonde sample every method here fails it, and the failure is information about the data.

\paragraph{Against DML.} The recommendation comes with the fold count attached: on the same data sets, DML with a cross validated lasso covers 0.694 with two folds against 0.826 with five at $s = 15$, and 0.450 against 0.734 at $s = 30$. We also say what would change our minds. A bootstrap that refits the nuisance, shown to estimate the dispersion of the bias across fits consistently, would reduce the structural argument for the posterior to convenience. The one we ran in Section~\ref{sec:crit:se} does not: it overstates the variance by a factor of 1.7 to 2.8 in the sparse designs, and where the error is bias it lifts coverage only from 0.536 to 0.607 under weak overlap and not at all in the dense design. A repeated split median that covered at the nominal level under weak overlap, rather than at 0.60, would remove the one column that favours the posterior on coverage. A cosine of order one between the two nuisance errors under the posterior in the sparse regime would show its good centring to be a coincidence of these designs. And a correctly specified design in which the posterior reports a calibration ratio far from one would defeat the case.

\section*{Notes}

The product form of the second order remainder, a bilinear form in two nuisance errors, is the mixed bias property of Rotnitzky, Smucler and Robins (2021). The horseshoe is from Carvalho, Polson and Scott (2010), its theory from van der Pas and coauthors (2014, 2017a, 2017b), and the R learner from Nie and Wager (2021).

\begin{exercises}
\item\label{ex:chcrit:1} For binary $T$ in the partially linear model with effect $\tau(X)$, show that the DML estimand weights $\tau(x)$ by $e(x)\{1 - e(x)\}$, and find a covariate law and effect function for which it differs from the average effect by more than the average effect itself.
\item\label{ex:chcrit:2} Show that the third term of the bound \eqref{eq:dml:bound} vanishes when $\tau(x)$ is constant, and compute $V_{\mathrm{eff}}$ when $e(x) = e$ is constant and $\sigma_1^2 = \sigma_0^2 = \sigma^2$.
\item\label{ex:chcrit:3} Construct an example in which the propensity and outcome models are both misspecified and the augmented estimator has larger bias than regression adjustment.
\item\label{ex:chcrit:4} Using the 401(k) partitions of Chapter~\ref{ch:k401}, compute the probability, across random splits, that the cross fitted lasso estimate falls outside the posterior 50\% interval of the coupled horseshoe.
\item\label{ex:chcrit:5} A funder adopts a program at cost $k$ and receives the benefit $\mathcal{U}(\tau) = \min(\tau, c)$ when the effect is $\tau$, so the benefit is linear in the effect up to the threshold $c$ and constant beyond it. Show that the Bayes action is to adopt when $\E\{\min(\tau, c) \mid \text{data}\} > k$, that this expected benefit equals $\E(\tau \mid \text{data}) - \E\{(\tau - c)_+ \mid \text{data}\}$, and that two posteriors with the same mean and variance can lead to opposite decisions.
\item\label{ex:chcrit:6} In the setting of Proposition~\ref{crit:prop:alignment} give the treatment and outcome regressions separate penalties $\lambda_{\mathrm L}^m$ and $\lambda_{\mathrm L}^\ell$. Show that to leading order $b = \lambda_{\mathrm L}^m s(\lambda_{\mathrm L}^\ell - \tau\lambda_{\mathrm L}^m)$, so that the cross channel cancels the own square channel exactly when $\lambda_{\mathrm L}^\ell = \tau\lambda_{\mathrm L}^m$. If each penalty is the same constant multiple of the noise standard deviation of its equation and $\tau > 0$, show that cancellation requires $\sigma_\varepsilon = 0$, and give the sign of $b$ at $\tau = 1$.
\item\label{ex:chcrit:7} Derive \eqref{eq:crit:cfvar} and show that if $\mathrm{Cov}(\bar\Delta_n, U_n) = 0$ the variance omitted by the sandwich is $\mathrm{Var}(\bar\Delta_n)$. In the least squares design of Section~\ref{sec:crit:se} at $K = 5$ the terms are $\mathrm{Var}(\hat\tau) = 3.540\times 10^{-3}$, $\mathrm{Var}(\bar\Delta_n) = 2.151\times 10^{-4}$, $2\,\mathrm{Cov}(\bar\Delta_n, U_n) = -2.365\times 10^{-4}$ and $\mathrm{Var}(U_n) = 3.562\times 10^{-3}$. Check the identity, compute the ratio of the cross term to $\mathrm{Var}(\bar\Delta_n)$, and say whether $\mathrm{Var}(U_n)$, the block the sandwich is meant to estimate, is too large or too small for $\mathrm{Var}(\hat\tau)$ here.
\item\label{ex:chcrit:8} The multiplier bootstrap at a fixed cross fit draws $\hat\tau^* = \hat\tau + \hat J_n^{-1}\,n^{-1}\sum_i \xi_i\hat\nu_i\hat\varepsilon_i$, where $\hat\nu_i = T_i - \hat m(X_i)$ and $\hat\varepsilon_i = Y_i - \hat\ell(X_i) - \hat\tau\hat\nu_i$ are the cross fitted residuals, $\hat J_n = n^{-1}\sum_i\hat\nu_i^2$, and the $\xi_i$ are independent with mean zero and variance one. Show that the variance of $\hat\tau^*$ given the data and the fits is the sandwich variance estimator, and explain why the term $B_n$ of Lemma~\ref{dml:lem:exact}, a constant given the data and the fits, moves the location of $\hat\tau$ but adds nothing to the variation of $\hat\tau^* - \hat\tau$ across bootstrap draws.
\item\label{ex:chcrit:9} For the Gaussian family of Proposition~\ref{crit:prop:dufour} show that $\mathrm{KL}(P_{j,\tau}^{\otimes n}\,\|\,P_{j,0}^{\otimes n}) = n\tau^2 j/(2\sigma_\varepsilon^2)$. Apply Pinsker's inequality and the triangle inequality through $P_{j,0}$ to obtain the bound $L\sqrt{nj}/\sigma_\varepsilon$ on $\|P_{j,\tau}^{\otimes n} - P_{j,\tau'}^{\otimes n}\|_{\mathrm{TV}}$ for $|\tau|, |\tau'| \le L$. For $n = 300$, $\sigma_\varepsilon = 1$ and $L = 5$, find the largest $j$ for which the bound is at most 0.1.
\item\label{ex:chcrit:10} Eight equal strata have propensities spaced evenly from 0.06 to 0.50 and stratum effects spaced evenly from 2.0 down to 0.0. Compute the weights $e_g(1 - e_g)$, their sum, the weighted sum of the effects and the ratio of the two, and compare it with the average effect. Which stratum gets the smallest weight and which the largest? Reverse the order of the effects, 0.0 at the lowest propensity and 2.0 at the highest, and recompute: does the partially linear score now overstate the average effect?
\end{exercises}
